\chapter{The big picture}
\label{ch:bigpicture}

A protein structure from the Protein Data Bank is a photograph, not a
simulation-ready model. It may be missing atoms, contain extra molecules from
the crystal, carry no hydrogen atoms, and include metal centers or cofactors
that the simulation software does not recognize. ProPrep takes you from that
photograph to a set of files that a simulation program can run, one decision at
a time.

This chapter shows the whole journey on one page, explains what each stage is
for, and tells you which chapters you need for your project.

\section{The stages of preparation}
\label{sec:bigpicture-stages}

\begin{figure}[ht]
\centering
\begin{tikzpicture}[
  node distance=5mm and 4mm,
  stage/.style={rectangle,rounded corners=3pt,draw=ppblue,fill=ppblue!7,line width=0.7pt,
    minimum width=37mm,minimum height=11mm,align=center,font=\sffamily\small},
  opt/.style={rectangle,rounded corners=3pt,draw=ppteal,dashed,fill=ppteal!6,line width=0.6pt,
    minimum width=31mm,minimum height=9mm,align=center,font=\sffamily\footnotesize},
  keyl/.style={font=\ttfamily\scriptsize\bfseries,text=ppteal},
  arr/.style={-{Stealth[length=2.4mm]},line width=0.8pt,color=ppblue},
  side/.style={-{Stealth[length=2mm]},line width=0.6pt,color=ppteal,dashed}]
  \node[stage] (load)  {\textbf{1. Get a structure}\\PDB, AlphaFold, own file};
  \node[stage,below=of load] (clean) {\textbf{2. Clean and repair}\\remove, fix, mutate};
  \node[stage,below=of clean] (prot) {\textbf{3. Protonation}\\charges at your pH};
  \node[stage,below=of prot] (special) {\textbf{4. Special parts}\\redox sites, metals, ligands};
  \node[stage,below=of special] (top) {\textbf{5. Build the system}\\force field, water, ions};
  \node[stage,below=of top] (md) {\textbf{6. Simulate}\\minimize, heat, run};
  \node[stage,below=of md] (ana) {\textbf{7. Analyze}\\RMSD, energies, distances};
  \foreach \a/\b in {load/clean,clean/prot,prot/special,special/top,top/md,md/ana}
    \draw[arr] (\a) -- (\b);
  \node[keyl,left=1mm of load]    {1a 1b};
  \node[keyl,left=1mm of clean]   {4a 4c 4b};
  \node[keyl,left=1mm of prot]    {4e 4f};
  \node[keyl,left=1mm of special] {2a 4d 2b 2c};
  \node[keyl,left=1mm of top]     {5a};
  \node[keyl,left=1mm of md]      {5b};
  \node[keyl,left=1mm of ana]     {5b};
  \node[opt,right=14mm of load]    (cmp)  {Compare structures\\\texttt{3a 3b}};
  \node[opt,right=14mm of clean]   (view) {3D viewer\\\texttt{ua}};
  \node[opt,right=14mm of special] (mem)  {Membrane builder\\\texttt{4g}};
  \node[opt,right=14mm of top]     (dock) {Docking\\\texttt{5d}};
  \node[opt,right=14mm of md]      (qm)   {QM/MM inputs\\\texttt{5c}};
  \draw[side] (load) -- (cmp);
  \draw[side] (clean) -- (view);
  \draw[side] (mem) -- (top);
  \draw[side] (prot.east) to[out=0,in=180] (dock.west);
  \draw[side] (top.east) to[out=0,in=180] (qm.west);
\end{tikzpicture}
\caption{The main route through ProPrep (left, solid) and the side tools you
use only when your project needs them (right, dashed). The codes are the keys
you type in ProPrep's main menu.}
\label{fig:route}
\end{figure}

\begin{description}[style=nextline,leftmargin=0pt,itemsep=1pt,parsep=0pt,font=\sffamily\bfseries\color{ppblue}]
\item[1. Get a structure (Chapter~\ref{ch:structure})]
  Download a structure from the Protein Data Bank or AlphaFold, or open a file
  you already have. Decide whether you want the molecule as it sits in the
  crystal or as it works in the cell.
\item[2. Clean and repair (Chapter~\ref{ch:cleaning})]
  Remove what you do not want to simulate (crystal water, salts, extra copies),
  find missing pieces, and change amino acids if your experiment needs a mutant.
\item[3. Protonation (Chapter~\ref{ch:protonation})]
  Decide which amino acids carry a charge at your pH. This sets the electric
  charges that drive most of the protein's behavior in a simulation.
\item[4. Special parts (Chapters~\ref{ch:redox} and~\ref{ch:forcefields})]
  Hemes, iron--sulfur clusters, metals, drugs and other non-standard molecules
  need extra parameters before a simulation can handle them.
\item[5. Build the system (Chapter~\ref{ch:topology})]
  Choose the force field, surround the protein with water and salt, and let
  ProPrep write the two files every Amber simulation needs: the
  \term{topology} (what is connected to what, and how strongly) and the
  \term{coordinates} (where every atom is).
\item[6. Simulate (Chapter~\ref{ch:runningmd})]
  Relax the structure, warm it to body temperature, and run the simulation on
  your computer or a computing cluster.
\item[7. Analyze (Chapter~\ref{ch:analysis})]
  Measure how the protein moved: how far it drifted from the start, which parts
  are flexible, and how energies changed.
\end{description}

You do not have to do every stage, and you can stop after any of them. Many
people use ProPrep only for stages 1 to 5 and run the simulation elsewhere.

\section{Which chapters do I need?}
\label{sec:bigpicture-which}

Find the row that matches your project. Read the chapters in the order given.

\begin{center}
\begin{tabularx}{\linewidth}{@{}>{\raggedright\arraybackslash}p{0.38\linewidth}L@{}}
\toprule
\textbf{Your project} & \textbf{Chapters, in order} \\
\midrule
First time with ProPrep & \ref{ch:starting}, \ref{ch:menus}, then the tutorial \ref{ch:tutorial} \\
An ordinary protein in water & \ref{ch:structure}, \ref{ch:cleaning}, \ref{ch:protonation},
  \ref{ch:topology}, \ref{ch:runningmd}, \ref{ch:analysis} \\
A protein with a heme, iron--sulfur cluster, flavin or copper site &
  \ref{ch:structure}, \ref{ch:cleaning}, \ref{ch:redox}, \ref{ch:protonation},
  \ref{ch:topology}, \ref{ch:runningmd} \\
A protein with a bound drug, ligand or unusual residue &
  \ref{ch:structure}, \ref{ch:cleaning}, \ref{ch:forcefields}, \ref{ch:protonation},
  \ref{ch:topology} \\
Finding where a small molecule binds & \ref{ch:structure}, \ref{ch:cleaning}, \ref{ch:docking} \\
A protein that sits in a cell membrane & \ref{ch:structure}, \ref{ch:cleaning},
  \ref{ch:protonation}, \ref{ch:membrane}, \ref{ch:topology}, \ref{ch:runningmd} \\
A QM/MM calculation on an active site & \ref{ch:redox}, \ref{ch:topology}, \ref{ch:qmmm} \\
Comparing related proteins or models & \ref{ch:structure}, \ref{ch:compare} \\
The same preparation for many proteins & do one by hand, then \ref{ch:sessions}, \ref{ch:batch} \\
\bottomrule
\end{tabularx}
\end{center}

\section{What ProPrep needs from you}
\label{sec:bigpicture-decisions}

ProPrep does the tedious work: downloading, renaming atoms, writing input files,
running the Amber programs. It does not make the scientific decisions for you.
At each step it shows you what it found and suggests a default, and you decide.
The decisions that matter most are:

\begin{itemize}
  \item \textbf{What to keep.} Which chains, cofactors, metals and water
        molecules belong in your model (Chapter~\ref{ch:cleaning}).
  \item \textbf{The pH.} It decides which residues are charged
        (Chapter~\ref{ch:protonation}).
  \item \textbf{Oxidation states.} For redox proteins, whether each site is
        oxidized or reduced (Chapter~\ref{ch:redox}).
  \item \textbf{The force field and water model.} The physical rules of the
        simulation (Chapter~\ref{ch:topology}).
  \item \textbf{How long to simulate.} ProPrep's protocols are starting points;
        the right length depends on your question (Chapter~\ref{ch:runningmd}).
\end{itemize}

\begin{tip}
ProPrep records every answer you give in a \term{session file}. If you change
your mind about an early decision, you do not have to start again: you can go
back and change that one answer (Chapter~\ref{ch:sessions}).
\end{tip}

\section{What works out of the box, and what needs more}
\label{sec:bigpicture-software}

Most of ProPrep works as soon as it is installed. A few tools call programs that
you must obtain separately. Appendix~\ref{app:software} lists every tool; in
short:

\begin{center}
\begin{tabularx}{\linewidth}{@{}lL@{}}
\toprule
\textbf{Extra program} & \textbf{Needed for} \\
\midrule
MODELLER license key (free for academics) & filling in missing residues and atoms,
  standard mutations, homology models (Chapters~\ref{ch:cleaning}, \ref{ch:compare}) \\
Gaussian (commercial) & new charges for ligands, modified residues and metal sites;
  running QM/MM inputs (Chapters~\ref{ch:forcefields}, \ref{ch:qmmm}) \\
ORCA (free for academics) & ORCA-style QM/MM inputs (Chapter~\ref{ch:qmmm}) \\
pmemd or a GPU (Amber license) & faster simulations; the free \texttt{sander}
  program works without them (Chapter~\ref{ch:runningmd}) \\
Internet connection & downloading structures, BLAST searches \\
\bottomrule
\end{tabularx}
\end{center}
